\documentclass[12pt,a4paper]{report}

\usepackage[utf8]{inputenc} 
\usepackage[T1]{fontenc}
\usepackage[french]{babel}
\usepackage{graphicx}
\usepackage[a4paper, margin=2.5cm]{geometry}
\usepackage{setspace}
\usepackage{float}
\usepackage{tikz}
\usetikzlibrary{arrows,positioning,shapes,decorations.pathreplacing}


\begin{document}

\begin{titlepage}
    \thispagestyle{empty}

    \begin{flushright}
        Année 2025--2026
    \end{flushright}

    \vspace{1.5cm}

    \begin{center}
        \includegraphics[width=1\textwidth]{../images/unige\_csd.png}\hspace{1cm}
    \end{center}

    \vspace{1.5cm}

    \begin{center}
        {\LARGE Projet de Bachelor d’informatique}
    \end{center}

    \vspace{0.5cm}

    \begin{center}
        {\Huge\bfseries SOM : Self-Organizing Maps}
    \end{center}

    \vfill

    \begin{center}
        {\Large William Guillermo Calero Roncal}
    \end{center}

    \vfill

    \begin{center}
        Année 2025--2026
    \end{center}

\end{titlepage}

% =========================
% Table des matières
% =========================
\tableofcontents
\clearpage

% =========================
% Contenu chap1
% =========================
\chapter{Introduction}

Durant les dernières décennies, plusieurs méthodes d'apprentissage automatique ont fait leur apparition. 
Cela s'est produit par une croissance exponentielle des données récoltées avec lesquelles nous tentons de répondre à 
des questions dans plusieurs domaines. Plusieurs méthodes s'inspirent directement du cerveau humain et une d'entre elles 
se nomme les maps auto-organisées (self organizing-maps: SOM). Cette méthode s'inspire directement des brain maps. 
Certaines cellules du cerveau répondent de manière sélective à un stimulus sensoriel spécifique et ces cellules ont tendance 
à s'assembler formant une topographie particulière en fonction du stimulus spécifique \cite{Kohonen2012SOM}.
Les SOM permettent la projection des données de grande dimension dans une grille à petite dimension tout en essayant de préserver leur topologie locale en tenant compte 
de leur voisinage.

\section{Contexte}
Actuellement, les réseaux de neurones ont une place centrale dans le domaine de l'apprentissage automatique. 
Cette méthode, s'inspirant aussi des neurones du cerveau humain, consiste à former des interconnexions entre 
couches des neurones artificiels. Dans la littérature, nous observons une tendance à appeler les SOM des réseaux 
de neurones. Cependant, nous observons certaines différences remarquables entre les réseaux de neurones et les maps 
auto-organisées. Pour rester au clair, nous allons rappeler la définition de ce qu'est un neurone artificiel, puis nous 
le comparerons avec le “neurone” d'un SOM.

\subsection{Neurone artificiel}
Un neurone artificiel ou plus communément un neurone formel est une représentation d'une combinaison linéaire suivie d'une fonction d'activation \cite{Cours_IA_UNIGE}. 
Pour un neurone $k$ nous avons $m+1$ entrées avec pour chaque entrée des données 
correspondantes $x_0,\dots, x_m$ et des poids $w_{k0}, \dots, w_{km}$. À cela nous ajoutons une fonction d'activation $\varphi$ qui
nous donne une valeur en sortie qui va être par la suite être transférée à la couche suivante (si ce n'est la dernière). 
De manière plus formel la donnée de sortie pour le neurone $k$ est donnée par l'équation suivante : 
$y_k = \varphi(\sum_{i=0}^{m} w_{ki} x_i)$ \cite{Artificial_neuron_wikipedia}.
Ce type de neurone utilise une fonction de coût (loss function) lors de la rétropropagation dont le but est de la minimiser 
l'erreur d'approximation. Dans ce contexte, nous avons un apprentissage supervisé, car la fonction se repose sur la connaissance des labels. 
Ainsi, un neurone artificiel est une analogie qui interprète  mathématiquement 
le fonctionnement d'un neurone avec des méthodes statistiques de machine learning (apprentissage automatique).
\begin{figure}[H]
\centering
\begin{tikzpicture}[
  neuron/.style={circle, draw, minimum size=1.5cm},
  input/.style={circle, draw, fill=green!20, minimum size=0.8cm},
  neu_som/.style={circle, draw, minimum size=2.3cm},
  >=latex
]   
    % neurone artificiel
    \node[input] (x0) at (0,3) {$x_0$};
    \node[input] (x1) at (0,1.5) {$x_1$};
    \node[input] (x2) at (0,0) {$x_2$};
    \node at (0,-1) {$\vdots$};
    \node[input] (xn) at (0,-2) {$x_n$};

    \node at (8,1) {$y = \varphi(z)$};
    
    \node[neuron] (n) at (4,1) {$z = \sum w_i x_i$};

    \draw[->] (x0) -- node[midway, above left] {$w_0$} (n);
    \draw[->] (x1) -- node[midway, above left] {$w_1$} (n);
    \draw[->] (x2) -- node[midway, above left] {$w_2$} (n);
    \draw[->] (xn) -- node[midway, above left] {$w_n$} (n);
    \draw[->] (n) -- (7,1);
\end{tikzpicture}
\caption{Neurone formel}
\label{fig:neurone}
\end{figure}

\subsection{Neurone du SOM}
Par contraste, un "neurone" d'un SOM est une unité de calcul qui représente un prototype de données.
Il est défini par un vecteur de poids $m$ de même dimension que les données d'entrée. 
Ce "neurone" est mis à jour en fonction de la distance entre les données d'entrée et le neurone lui-même. 
Le neurone qui est le plus proche de la donnée d'entrée est appelé Best Matching Unit (BMU) et c'est lui qui est mis à jour
ainsi que ses voisins en fonction de la formule suivante : $m_c(t+1) = m_c(t) + h_{ci}(t)[\tilde{x} - m_c(t)]$ \cite{Kohonen2012SOM}. 
Dans ce contexte, nous avons un apprentissage non supervisé, car nous n'avons pas de labels pour guider l'apprentissage.
Cela est une différence fondamentale avec les neurones formels, car les SOM sont utilisés pour la visualisation et la réduction de la dimensionnalité, 
tandis que les neurones formels sont utilisés plutôt pour la classification et la régression. 
De plus, les SOM préservent la topologie des données d'entrée, ce qui n'est pas le cas des neurones formels.

Même s'ils partagent une inspiration biologique commune, les SOM présentent des caractéristiques uniques qui les distinguent des réseaux de neurones traditionnels.
Ainsi, même si historiquement les maps auto-organisées sont classées comme des réseaux de neurones, par les différences que nous avons citées,
il est raisonnable de les considérer comme des structures différentes. 
Par la suite de ce rapport et pour éviter toutes confusions, nous allons utiliser "noeud" pour les unités de calcul d'un SOM.

\begin{figure}[H]
\centering
\begin{tikzpicture}[
  neuron/.style={circle, draw, minimum size=1.5cm},
  input/.style={circle, draw, fill=green!20, minimum size=0.8cm},
  neu_som/.style={circle, draw, minimum size=2.3cm},
  >=latex
]   
    
    % neurone d'un SOM
    \node[neu_som] (x) at (0,-6) {$\tilde{x}$};
    
    \node[neu_som] (mc) at (5,-6) {$m_c$};
    
    \draw[->] (x) -- node[midway, above] {$BMU$} (mc);
    \draw[->, thick, loop above] (mc) to [out=0, in=120, looseness=5] node[above, rotate=-30] {$m_c + h_{ci}[\tilde{x} - m_c] $} (mc);
\end{tikzpicture}
\caption{Noeud d'un SOM}
\label{fig:neoudSOM}
\end{figure}

\section{Objectifs du projet}

L'objectif de ce projet consiste à comprendre les maps auto-organisées, leurs limites ainsi que 
les extensions proposées dans la littérature pour y remédier. Dans la suite de ce rapport, un état de l'art sur les SOM sera présenté. 
Ensuite, nous ferons une étude empirique en prenant un SOM from scratch, un SOM d'une bibliothèque existante ainsi qu'une l'extension D-SOM. 
Enfin, nous analyserons et discuterons de leurs performances, afin de mettre en évidence leurs avantages et inconvénients.


% =========================
% Contenu chap2
% =========================

\chapter{État de l’art}

\section{Self-Organizing Maps (SOM)}

\section{Limites et Extensions}

\section{D-SOM}


% =========================
% Contenu chap3
% =========================

\chapter{Etude empirique}

\section{Données utilisées}

\section{SOM de base}

\section{D-SOM}

\section{Comparaison et analyse des résultats}

% =========================
% Contenu chap4
% =========================

\chapter{Conclusion}



\bibliographystyle{plain}
\bibliography{bibliography}

\end{document}